\section*{Conclusão e Diretrizes Operacionais}
\addcontentsline{toc}{section}{Conclusão e Diretrizes Operacionais}

A revisão bibliográfica demonstra que a graduação em escalada não é um juízo arbitrário, mas um sistema multidimensional controlável. Cinco diretrizes operacionais sintetizam a aplicação prática para qualquer ginásio brasileiro:

\begin{enumerate}
    \item \textbf{Vias-Gabarito (\textit{Benchmark Boulders})}: Manter de 8 a 12 blocos de referência calibrados com polos de rocha nacional (Cocalzinho e Igatu) para ancorar os circuitos.
    \item \textbf{Forerunning Obrigatório}: Proibir a graduação por um único \textit{setter}, exigindo validação por testadores de estaturas distintas para eliminar viés morfológico e \textit{beta breaks} \cite{cwa_setting, anderson_routesetting}.
    \item \textbf{Proteção da Base}: Separar a iniciação absoluta (\texttt{VB/V0}) da transição (\texttt{V0/V1}) por meio de faixas distintas ou sub-identificadores claros na base da via \cite{ircra_draper, cwa_setting}.
    \item \textbf{Governança Híbrida}: Combinar parede física limpa (fitas/tags) com registro unitário em aplicativo digital, auditando divergências superiores a 1 grau após 15 repetições \cite{topclimbing_platform, pmc_ric_study}.
    \item \textbf{Vínculo Institucional}: Cadastrar o ginásio no SGE da CBEscalada para habilitar atletas vinculados ao calendário oficial e programas de fomento \cite{cbescalada_filiacoes}.
\end{enumerate}
